% P15-U1: full measure-extension and distribution-existence proofs.
\section{Constructing measures and distributions: proofs}
\label{sec:annex-p15-measure-extension}

These arguments use Chapter~\ref{sec:p15-measurable-probability} for
nonnegative sums, continuity of measures, Borel generators, pushforward
laws and the pi-lambda theorem. The definitions of the present chapter
are in Chapter~\ref{sec:p15-measure-extension}. The source is Durrett's
11 January 2019 author version
\cite[thms.~1.1.4, 1.1.9 and 1.2.2; app.~A.1]{DurrettPTEAuthor2019}.

\begin{annexitem}{annex:p15.semialgebra-extension}{thm:p15-semialgebra-extension}{From elementary pieces to a premeasure}
\begin{proof}
Write \(\mathcal A=\mathcal A(\mathcal S)\). The empty union puts
\(\varnothing\) in \(\mathcal A\), and the complement property
applied to \(\varnothing\in\mathcal S\) puts \(\Omega\) in it.
If \(A=\bigcup_{i=1}^r S_i\) and \(B=\bigcup_{j=1}^s T_j\)
are disjoint representations, then the intersections \(S_i\cap T_j\)
are members of \(\mathcal S\), are pairwise disjoint, and have union
\(A\cap B\). Thus \(\mathcal A\) is closed under binary
intersections. For each \(S_i\), its complement belongs to
\(\mathcal A\) by the semialgebra property. Consequently
\(A^c=\bigcap_{i=1}^r S_i^c\) belongs to \(\mathcal A\) by
finite intersection closure, including \(r=0\), when it equals
\(\Omega\). Complements and binary intersections give binary unions,
so \(\mathcal A\) is an algebra. Any algebra containing \(\mathcal S\)
contains all its finite unions and hence \(\mathcal A\), proving
minimality. A sigma-algebra contains \(\mathcal S\) exactly when it
contains \(\mathcal A\), which proves equality of their generated
sigma-algebras.

Suppose \(m\) has the stated finite partition property. If
\(A=\bigcup_i S_i=\bigcup_j T_j\) are two disjoint finite
representations, each \(S_i\) is partitioned by the sets
\(S_i\cap T_j\), and each \(T_j\) by the sets \(S_i\cap T_j\).
The partition property and finite regrouping of nonnegative terms give
\[
 \sum_i m(S_i)=\sum_i\sum_j m(S_i\cap T_j)
             =\sum_j\sum_i m(S_i\cap T_j)=\sum_j m(T_j).
\]
Thus \(\overline m\) is well defined. If finitely many algebra
members \(A_1,\ldots,A_k\) are disjoint, the elementary pieces in
their representations form one disjoint finite representation of their
union. This proves finite additivity of \(\overline m\). Any finitely
additive extension must have the displayed sum on a representation, so
it is unique. Its value at the empty set is zero.

For \(A\subseteq B\) in the algebra, finite additivity gives
\(\overline m(B)=\overline m(A)+\overline m(B\setminus A)\),
proving monotonicity without subtraction. For a finite cover of \(A\),
put \(C_1=B_1\) and
\(C_j=B_j\setminus\bigcup_{i<j}B_i\) for \(j>1\). These are
disjoint algebra members, \(C_j\subseteq B_j\), and they cover
\(A\). Their intersections with \(A\) partition \(A\). Hence
\[
 \overline m(A)=\sum_{j=1}^k\overline m(A\cap C_j)
       \leq\sum_{j=1}^k\overline m(B_j).
\]
For \(k=0\), the cover condition forces \(A=\varnothing\), and
the inequality is equality of two zeros.

Now assume the additional countable partition inequality on \(\mathcal S\).
Let \(A=\bigcup_{n\geq1}A_n\) be a disjoint union with all these
sets in \(\mathcal A\). Represent each \(A_n\) by finitely many
elementary pieces. Taken together these pieces form a countable disjoint
family \((T_k)\) in \(\mathcal S\) with union \(A\); finite or
empty families can be padded by empty sets. The nonnegative grouping
theorem gives
\(\sum_n\overline m(A_n)=\sum_k m(T_k)\).
Choose a finite disjoint representation \(A=\bigcup_{j=1}^r S_j\).
For each \(j\), the sets \(S_j\cap T_k\) partition \(S_j\), so
\[
 \begin{aligned}
 \overline m(A)=\sum_{j=1}^r m(S_j)
  &\leq\sum_{j=1}^r\sum_{k\geq1}m(S_j\cap T_k)\\
  &=\sum_{k\geq1}\sum_{j=1}^r m(S_j\cap T_k)
   =\sum_{k\geq1}m(T_k)=\sum_{n\geq1}\overline m(A_n).
 \end{aligned}
\]
The last finite sums equal \(m(T_k)\) by the finite partition property,
since \(T_k\subseteq A\). Regrouping is valid even when a sum is
infinite. Conversely, finite additivity and monotonicity give
\(\sum_{n=1}^N\overline m(A_n)\leq\overline m(A)\) for every
\(N\). Taking the supremum of these partial sums gives the reverse
inequality. This is precisely the premeasure property.
\end{proof}
\end{annexitem}

\begin{annexitem}{annex:p15.caratheodory-measurable}{thm:p15-caratheodory-measurable}{The measurable domain of an outer measure}
\begin{proof}
For any \(E,T\subseteq\Omega\), countable subadditivity, with two
terms and empty sets thereafter, gives
\(q(T)\leq q(T\cap E)+q(T\setminus E)\). Thus only the reverse
inequality must be proved to establish the splitting criterion.
Both \(\varnothing\) and \(\Omega\) satisfy it, and the criterion
is symmetric between \(E\) and \(E^c\).

Suppose \(E,F\in\mathcal M_q\). Splitting first by \(E\) and
then splitting \(T\setminus E\) by \(F\) gives
\[
 q(T)=q(T\cap E)+q((T\setminus E)\cap F)
                         +q(T\setminus(E\cup F)).
\]
The first two sets have union \(T\cap(E\cup F)\). Their outer
measures have sum at least \(q(T\cap(E\cup F))\) by subadditivity.
Hence \(E\cup F\) satisfies the criterion. Complementation now
gives closure under intersections and differences, so \(\mathcal M_q\)
is an algebra. The displayed equalities use only nonnegative addition,
including when values are infinite.

For disjoint \(E_1,\ldots,E_n\in\mathcal M_q\) and any \(T\),
repeated splitting gives
\[
 q\left(T\cap\bigcup_{i=1}^n E_i\right)
                   =\sum_{i=1}^n q(T\cap E_i).
\]
To see the induction step explicitly, split
\(T\cap\bigcup_{i=1}^n E_i\) by \(E_n\); disjointness makes
the two pieces \(T\cap E_n\) and
\(T\cap\bigcup_{i=1}^{n-1}E_i\). The case \(n=0\) is zero.

For a disjoint sequence \((E_i)\) in \(\mathcal M_q\), let
\(E=\bigcup_{i\geq1}E_i\) and \(F_n=\bigcup_{i=1}^n E_i\).
Since \(F_n\in\mathcal M_q\) and \(T\setminus E\subseteq
T\setminus F_n\), splitting and monotonicity imply
\[
 q(T)\geq\sum_{i=1}^n q(T\cap E_i)+q(T\setminus E).
\]
Taking the supremum over \(n\), addition of a fixed nonnegative
extended value commutes with this increasing supremum. Indeed, if that
fixed value is infinite both sides of the assertion are infinite; if
it is finite, the assertion follows from ordinary addition and the
definition of an infinite increasing limit. We obtain
\[
 q(T)\geq\sum_{i\geq1}q(T\cap E_i)+q(T\setminus E)
       \geq q(T\cap E)+q(T\setminus E).
\]
The second inequality is subadditivity. Thus \(E\in\mathcal M_q\).
An arbitrary sequence of members of \(\mathcal M_q\) can be made
disjoint by replacing each set by its difference with the union of its
predecessors. Those differences belong to its algebra, preserve the
union, and the preceding argument applies. Consequently
\(\mathcal M_q\) is a sigma-algebra.

For the disjoint sequence just used, monotonicity and finite splitting
with \(T=\Omega\) give
\(q(E)\geq q(F_n)=\sum_{i=1}^n q(E_i)\). Supremizing proves
\(q(E)\geq\sum_i q(E_i)\). Subadditivity gives the other inequality,
so the restriction is countably additive and is a measure.

Finally, if \(q(N)=0\) and \(E\subseteq N\), monotonicity gives
\(q(E)=0\) and \(q(T\cap E)=0\) for every \(T\). Subadditivity
and monotonicity yield
\[
 q(T)\leq q(T\setminus E)+q(T\cap E)
       =q(T\setminus E)\leq q(T).
\]
This proves the splitting criterion for \(E\), and its restricted
measure is zero. It proves both the zero-outer-measure assertion and
completeness of the resulting measure.
\end{proof}
\end{annexitem}

\begin{annexitem}{annex:p15.premeasure-extension}{thm:p15-premeasure-extension}{Existence by covers and uniqueness by finite restrictions}
\begin{proof}
A premeasure is finitely additive, by appending empty sets, and monotone,
by decomposing a larger algebra member into the smaller one and its
difference. We first need its countable covering inequality. If
\(A,A_n\in\mathcal A\) and \(A\subseteq\bigcup_n A_n\), set
\[
 D_n=(A\cap A_n)\setminus\bigcup_{j<n}A_j.
\]
These disjoint algebra members have union exactly \(A\), which is
in the algebra even when \(\bigcup_n A_n\) is not. Thus the
premeasure property and monotonicity give
\[
 \mu_0(A)=\sum_n\mu_0(D_n)\leq\sum_n\mu_0(A_n).
\]

The empty set has a cover consisting of empty sets, of cost zero, so
\(\mu^*(\varnothing)=0\). Every cover of \(F\) also covers
\(E\subseteq F\), which proves monotonicity of the infimum.
To prove subadditivity, suppose first that
\(\sum_j\mu^*(E_j)<\infty\). Given \(\varepsilon>0\), choose
for each \(j\geq1\) a cover \((A_{j,n})_{n\geq1}\) of \(E_j\)
from the algebra with
\[
 \sum_n\mu_0(A_{j,n})<\mu^*(E_j)+\varepsilon2^{-j}.
\]
Such a cover exists because the infimum is finite. The doubly indexed
family is countable and covers \(\bigcup_j E_j\). Nonnegative
regrouping therefore gives
\[
 \mu^*\left(\bigcup_j E_j\right)
   \leq\sum_j\sum_n\mu_0(A_{j,n})
   \leq\sum_j\mu^*(E_j)+\varepsilon.
\]
The left side is finite by this bound, and letting \(\varepsilon\)
decrease to zero proves the desired inequality. If
\(\sum_j\mu^*(E_j)=\infty\), the desired inequality already holds
in the extended order. This proves that \(\mu^*\) is an outer measure.

For \(A\in\mathcal A\), the one-set cover followed by empty sets
gives \(\mu^*(A)\leq\mu_0(A)\). The premeasure covering inequality
proved above gives \(\mu_0(A)\leq\sum_n\mu_0(A_n)\) for every
cover of \(A\); taking their infimum gives the reverse inequality.
Hence \(\mu^*(A)=\mu_0(A)\), including infinite values.

Fix \(A\in\mathcal A\) and an arbitrary \(T\subseteq\Omega\).
If \(\mu^*(T)=\infty\), the inequality
\(\mu^*(T)\geq\mu^*(T\cap A)+\mu^*(T\setminus A)\) is automatic.
If \(\mu^*(T)<\infty\), choose a cover \((B_n)\) with total
cost less than \(\mu^*(T)+\varepsilon\). Finite additivity on
the algebra gives
\(\mu_0(B_n)=\mu_0(B_n\cap A)+\mu_0(B_n\setminus A)\).
The two sequences of pieces cover respectively \(T\cap A\) and
\(T\setminus A\), so
\[
 \begin{aligned}
 \mu^*(T)+\varepsilon
   &\geq\sum_n\mu_0(B_n\cap A)+\sum_n\mu_0(B_n\setminus A)\\
   &\geq\mu^*(T\cap A)+\mu^*(T\setminus A).
 \end{aligned}
\]
The sum on the last line is finite by this bound. Letting
\(\varepsilon\) tend to zero proves the reverse splitting inequality.
The forward inequality is outer subadditivity. Thus
\(A\in\mathcal M_{\mu^*}\).

The preceding outer-measure theorem now provides a measure on
\(\mathcal M_{\mu^*}\). Since this is a sigma-algebra containing
\(\mathcal A\), it contains \(\sigma(\mathcal A)\). Restriction
to that sigma-algebra gives the asserted extension. No part of this
existence proof used sigma-finiteness.

For uniqueness, let \(\nu_1,\nu_2\) be two measure extensions on
\(\sigma(\mathcal A)\), and let \(E_n\in\mathcal A\) be a
sigma-finite cover for \(\mu_0\). Put
\(H_n=\bigcup_{j=1}^n E_j\). These algebra members increase to
\(\Omega\) and have finite premeasure: disjointifying this finite
union gives \(\mu_0(H_n)\leq\sum_{j=1}^n\mu_0(E_j)<\infty\).
Thus \(\nu_1(H_n)=\nu_2(H_n)=\mu_0(H_n)<\infty\).
For fixed \(n\), define
\[
 \mathcal L_n=\{B\in\sigma(\mathcal A):
                  \nu_1(B\cap H_n)=\nu_2(B\cap H_n)\}.
\]
It contains \(\Omega\). If \(B\subseteq C\) are in
\(\mathcal L_n\), the measures of their intersections with \(H_n\)
are finite, and
\[
 \nu_i((C\setminus B)\cap H_n)
     =\nu_i(C\cap H_n)-\nu_i(B\cap H_n)
                  \qquad(i=1,2).
\]
The right sides agree, proving closure under nested differences.
For an increasing sequence \(B_k\in\mathcal L_n\), continuity
from below of each \(\nu_i\) proves equality on
\((\bigcup_k B_k)\cap H_n\). Thus \(\mathcal L_n\) is a
lambda-system. It contains \(\mathcal A\), because
\(B\cap H_n\in\mathcal A\) when \(B\in\mathcal A\).
The algebra is a pi-system. The proved pi-lambda theorem, the first
assertion of Theorem~\ref{thm:p15-measure-uniqueness}, gives
\(\sigma(\mathcal A)\subseteq\mathcal L_n\).
For any \(B\in\sigma(\mathcal A)\), we now have equality on
every \(B\cap H_n\). These sets increase to \(B\), so another
application of continuity from below proves \(\nu_1(B)=\nu_2(B)\),
whether this common value is finite or infinite. A finite premeasure
is sigma-finite by the cover consisting of \(\Omega\) itself, and
every extension preserves its total mass since \(\Omega\in\mathcal A\).
\end{proof}
\end{annexitem}

\begin{annexitem}{annex:p15.stieltjes-construction}{thm:p15-stieltjes-construction}{Stieltjes interval masses and their countable covering bound}
\begin{proof}
Intersections of the intervals \(I(a,b)\) are either empty or another
such interval, with left endpoint the larger left endpoint and right
endpoint the smaller right endpoint. The complement of \(I(a,b)\)
is the disjoint union of \(I(-\infty,a)\), if \(a\) is finite,
and \(I(b,\infty)\), if \(b\) is finite. Missing pieces are empty;
the complement of the empty set is \(I(-\infty,\infty)\).
This verifies the semialgebra axioms. All its intervals are Borel,
using closed rays and their complements. Conversely it contains every
closed ray \((-\infty,b]\), which generates the Borel sigma-algebra
by Theorem~\ref{thm:p15-borel-generators}. Thus
\(\sigma(\mathcal S_{\mathbb R})=\mathcal B(\mathbb R)\).

Because \(F\) is nondecreasing and finite at every real argument,
\(F(-\infty)\leq F(0)\leq F(\infty)\), with the endpoint
restrictions stated in the definition. Monotonicity makes these infimum
and supremum values its limits at the corresponding ends. For example,
if the infimum is finite, choose a point with value within
\(\varepsilon\) of it; all points to its left have values between
that infimum and the chosen value. If the infimum is \(-\infty\),
choose a point below any prescribed finite bound and use the same
monotonicity. The upper-end proof is the reversed argument.

In a finite disjoint partition of a nonempty interval into nonempty
members of \(\mathcal S_{\mathbb R}\), order the pieces from left
to right. The first left endpoint and last right endpoint equal the
target endpoints. Every successive left endpoint equals the preceding
right endpoint: a larger value would leave an uncovered point between
them, while a smaller value would create overlap. The masses therefore
add as consecutive increments of \(F\). If both extreme \(F\)
values are finite, ordinary telescoping gives the target increment.
If an extreme \(F\) value is infinite, the corresponding first or
last increment is infinite, and both the sum and the target mass are
infinite. This includes the single-piece case; all internal endpoints
of a partition with more pieces are finite. Empty pieces contribute
zero. Thus the finite partition property holds in every case.
The finite part of Theorem~\ref{thm:p15-semialgebra-extension} now
provides a finitely additive extension \(\overline m_F\) to the
interval algebra and its finite covering inequality. This use precedes
verification of the countable condition and needs only finite additivity.

We supply the real covering fact used next. Every cover of a closed
bounded interval \([\alpha,\beta]\) by open intervals has a finite
subcover. The case \(\alpha=\beta\) follows from one covering
interval. When \(\alpha<\beta\), let \(D\) consist of the
\(t\in[\alpha,\beta]\) for which \([\alpha,t]\) has a finite
subcover. It is nonempty because \(\alpha\in D\). Let
\(s=\sup D\). If \(s<\beta\), an open interval in the cover
containing \(s\) contains \((s-\delta,s+\delta)\) for some
\(\delta>0\). Choose \(t\in D\) with \(t>s-\delta\).
A finite cover of \([\alpha,t]\), together with that interval,
covers \([\alpha,r]\) for some \(s<r\leq\beta\) with
\(r<s+\delta\). This contradicts the definition of the supremum.
Hence \(s=\beta\). Choose a covering interval about \(\beta\)
and \(t\in D\) sufficiently close to \(\beta\) from below;
if \(t=\beta\), its finite cover already suffices. Otherwise that
finite cover and the interval about \(\beta\) cover the whole of
\([\alpha,\beta]\). This proves the covering fact using real
order completeness.

Consider a countable cover of a bounded target \((a,b]\), with
real \(a<b\), by intervals \(I(a_i,b_i)\). Choose finite
\(L<a<b<R\). Intersect every covering interval with \((L,R]\)
and call the result \(J_i\). The nonempty \(J_i\) have finite
endpoints, still cover \((a,b]\), and satisfy
\(m_F(J_i)\leq m_F(I(a_i,b_i))\) by the proved finite monotonicity.
This justifies reducing the covering pieces to bounded ones; it does
not assume their endpoints were finite originally.

Fix \(\varepsilon>0\). Right continuity of \(F\) at \(a\)
permits a \(\delta\) with \(0<\delta<b-a\) and
\(F(a+\delta)-F(a)<\varepsilon\). Write each nonempty
\(J_i=(\alpha_i,\beta_i]\). Right continuity at \(\beta_i\)
permits \(\eta_i>0\) such that
\[
 F(\beta_i+\eta_i)-F(\beta_i)<\varepsilon2^{-i}.
\]
The open intervals \((\alpha_i,\beta_i+\eta_i)\) cover
\([a+\delta,b]\), because each original piece is contained in its
open enlargement. Select a finite subcover with index set \(K\).
The corresponding half-open enlargements also cover \((a+\delta,b]\).
Applying the finite covering inequality in the interval algebra gives
\[
 \begin{aligned}
 F(b)-F(a+\delta)
  &\leq\sum_{i\in K}\bigl(F(\beta_i+\eta_i)-F(\alpha_i)\bigr)\\
  &\leq\sum_{i\geq1}m_F(J_i)+\varepsilon
   \leq\sum_{i\geq1}m_F(I(a_i,b_i))+\varepsilon.
 \end{aligned}
\]
Empty \(J_i\) contribute zero. Adding the finite error at the left
endpoint yields
\[
 F(b)-F(a)\leq\sum_{i\geq1}m_F(I(a_i,b_i))+2\varepsilon.
\]
If the series on the right is infinite, the desired covering inequality
is already true. If it is finite, letting \(\varepsilon\) decrease
to zero proves that inequality for bounded targets.

For a target \(I(a,b)\) with any allowed endpoints, its truncations
\(K_n=I(a,b)\cap(-n,n]\) are empty or bounded half-open intervals,
increase, and have union \(I(a,b)\). We verify directly from the
endpoint definitions, without measure continuity, that
\[
 m_F(I(a,b))=\sup_n m_F(K_n).
\]
The claim is immediate for a bounded target, since eventually
\(K_n=I(a,b)\). In general the nonempty truncations have endpoints
\(\max(a,-n)\) and \(\min(b,n)\). Their upper \(F\) values
increase to \(F(b)\), and their lower \(F\) values decrease to
\(F(a)\), by the endpoint-limit argument. If both limits are finite,
their differences converge to \(F(b)-F(a)\). If either is infinite,
choose a fixed real point \(c\) strictly between \(a\) and \(b\).
For all large \(n\), the increment of the truncation is at least each
of the nonnegative increments from its left endpoint to \(c\) and
from \(c\) to its right endpoint. At least one of these increments
tends to infinity. Thus the truncation masses increase to infinity,
as the endpoint assignment requires. This proves the supremum identity
in all cases with no subtraction of equal infinities.

Any given cover of \(I(a,b)\) also covers every nonempty \(K_n\).
The bounded-target inequality bounds each \(m_F(K_n)\) by the same
total cost of that cover. Taking the supremum gives the claimed general
countable covering bound. A disjoint partition is such a cover, so the
second semialgebra condition is verified.

The semialgebra theorem now gives a premeasure on its algebra. Its
sets \((-n,n]\), \(n\geq1\), cover \(\mathbb R\) and have
finite masses \(F(n)-F(-n)\), so it is sigma-finite. The premeasure
extension theorem gives a unique measure on
\(\sigma(\mathcal S_{\mathbb R})=\mathcal B(\mathbb R)\).
It preserves every elementary interval mass, including
\(I(-\infty,\infty)=\mathbb R\), proving the total-mass formula.
Any other Borel measure with those interval masses agrees on the algebra
by finite additivity and is the same measure by the proved uniqueness
argument. The resulting measure is itself sigma-finite by the same
bounded-interval cover.
\end{proof}
\end{annexitem}

\begin{annexitem}{annex:p15.lebesgue-unit-interval}{thm:p15-lebesgue-unit-interval}{Length measure and its probability restriction}
\begin{proof}
The function \(F(x)=x\) is real valued, nondecreasing and right
continuous. Its Stieltjes construction therefore gives a Borel measure
\(\lambda\) with \(\lambda((a,b])=b-a\) for real \(a<b\).
Its sigma-finite cover consists of the bounded intervals \((-n,n]\),
with finite measures \(2n\); their increasing union is \(\mathbb R\),
whose measure is infinity.

For a real \(c\), the intervals \((c-1/n,c]\) decrease to
\(\{c\}\) and have measures \(1/n\). Their first measure is
finite, so continuity from above gives \(\lambda(\{c\})=0\).
Changing the inclusion of either endpoint of a bounded interval adds
or removes only singleton null sets. Finite additivity therefore gives
the usual length in each case, including length zero for singletons
and empty intervals.

If any Borel measure has these bounded half-open interval lengths,
continuity from below applied to truncations by \((-n,n]\) determines
its value on every unbounded elementary interval as well. These are
the values assigned by \(m_F\) for \(F(x)=x\). It consequently
agrees on the entire interval algebra, so the Stieltjes uniqueness
theorem proves uniqueness from the stated bounded-interval data alone.

For real \(t\), translation and its inverse are continuous; hence
\(B+t\) is Borel whenever \(B\) is Borel, since it is the inverse
image of \(B\) under \(x\mapsto x-t\). The assignment
\(\lambda_t(B)=\lambda(B+t)\) is a measure: translation preserves
the empty set, disjointness and countable unions. It has the same
bounded half-open interval lengths, because
\((a,b]+t=(a+t,b+t]\) has length \(b-a\). The uniqueness just
proved gives \(\lambda_t=\lambda\), establishing translation
invariance.

The trace collection \(\mathcal F_U\) contains \((0,1)\), and is
closed under relative complements and countable unions because
\[
 (0,1)\setminus(B\cap(0,1))=B^c\cap(0,1),\qquad
 \bigcup_n(B_n\cap(0,1))=\left(\bigcup_nB_n\right)\cap(0,1).
\]
Each trace member is itself Borel in \(\mathbb R\), since \((0,1)\)
is Borel. Thus \(\mathbb U\) is defined there and inherits countable
additivity from \(\lambda\). Its total value is
\[
 \lambda((0,1))=\lambda((0,1])-\lambda(\{1\})=1-0=1.
\]
The subtraction is between finite numbers. The inverse image under
\(U(u)=u\) of any Borel set \(B\) is \(B\cap(0,1)\), so
\(U\) is measurable. Its sublevel event is empty when \(t\leq0\),
is \((0,t]\) when \(0<t<1\), and is all of \((0,1)\) when
\(t\geq1\). Their measures are respectively zero, \(t\) and one,
giving the stated cumulative distribution including both endpoints.
\end{proof}
\end{annexitem}

\begin{annexitem}{annex:p15.cdf-existence}{thm:p15-cdf-existence}{A prescribed cumulative function and its measurable quantile}
\begin{proof}
The hypotheses on \(F\) make it a Stieltjes measure function with
\(F(-\infty)=0\) and \(F(\infty)=1\). The measure already
constructed from its interval increments has total mass one and
\[
 \mu_F((-\infty,x])=F(x)-F(-\infty)=F(x).
\]
This proves existence of its Borel probability law before making the
quantile construction.

Fix \(u\in(0,1)\). The lower tail limit provides a real \(a\)
with \(F(a)<u\), so \(L_u=\{y:F(y)<u\}\) is nonempty. The
upper tail limit provides a real \(b\) with \(F(b)>u\).
Monotonicity then gives \(F(y)>u\) for all \(y\geq b\), so
\(b\) is an upper bound on \(L_u\). By real completeness its
supremum \(Q_F(u)\) is a finite real number. If \(u\leq v\),
then \(L_u\subseteq L_v\), so \(Q_F(u)\leq Q_F(v)\).

Let \(x\in\mathbb R\). If \(u\leq F(x)\), any
\(y\in L_u\) must satisfy \(y<x\): otherwise monotonicity
would give \(F(y)\geq F(x)\geq u\). Thus \(x\) is an upper
bound for \(L_u\) and \(Q_F(u)\leq x\).
If \(u>F(x)\), right continuity provides \(\delta>0\) with
\(F(x+\delta)<u\). For example, take the positive gap
\(u-F(x)\), use half that gap in the definition of right continuity,
and choose a sufficiently small positive \(\delta\).
Then \(x+\delta\in L_u\), whence
\(Q_F(u)\geq x+\delta>x\). These two implications prove the
ray equivalence for every \(x\), without assuming that \(F\) is
strictly increasing or that its left limits equal its values.

It follows that
\[
 \{u\in(0,1):Q_F(u)\leq x\}
                  =(0,1)\cap(-\infty,F(x)].
\]
The set on the right belongs to \(\mathcal F_U\). The real-ray
generator criterion makes \(Q_F\) measurable. If \(F(x)=0\),
the displayed set is empty; if \(F(x)=1\), it is \((0,1)\).
If \(0<F(x)<1\), it is \((0,F(x)]\), whose uniform measure
is its length \(F(x)\). Thus in all three cases
\(\mathbb U(Q_F\leq x)=F(x)\). The law of this random variable
is a Borel probability measure by the pushforward theorem. Its cumulative
function agrees with that of \(\mu_F\), so
Theorem~\ref{thm:p15-cdf-laws} proves equality of the two laws and
uniqueness among all Borel probability laws with cumulative function \(F\).

For fixed \(u\), the ray equivalence also gives the exact set identity
\[
 \{x\in\mathbb R:F(x)\geq u\}=[Q_F(u),\infty).
\]
In particular it includes its left endpoint; this uses right continuity
through the already proved equivalence. Its minimum is \(Q_F(u)\),
which proves the second description of the quantile and
\(u\leq F(Q_F(u))\). For every \(y<Q_F(u)\), the equivalence
gives \(F(y)<u\). Taking their supremum gives
\(F(Q_F(u)-)\leq u\). No assertion of strict inequality at that
left limit is needed.

To verify the jump-fiber bounds, suppose
\(F(x-)<u\leq F(x)\) and \(0<u<1\). The upper inequality
gives \(Q_F(u)\leq x\). For every \(y<x\),
\(F(y)\leq F(x-)<u\), so all such \(y\) belong to \(L_u\).
Their supremum is \(x\), giving \(Q_F(u)\geq x\). Hence
\(Q_F(u)=x\). Conversely, if \(Q_F(u)=x\), the quantile bounds
already proved give \(F(x-)\leq u\leq F(x)\). This proves both
inclusions, including the cases when either bound is zero or one.

The bounding intervals have uniform measure \(F(x)-F(x-)\).
Indeed their endpoints belong to \([0,1]\); intersecting with
\((0,1)\) can remove only the endpoints zero and one, and changing
the inclusion of \(F(x-)\) changes only one singleton. All those
singletons have Lebesgue measure zero by the previous theorem.
Monotonicity therefore gives the same measure to the intermediate
fiber. This also follows from the jump formula for its now established
law. Finally, if \(F(a)=F(b)\) with \(a<b\), the ray equivalence
makes \(\{Q_F\leq a\}=\{Q_F\leq b\}\) as sets. Their
difference \(\{a<Q_F\leq b\}\) is empty, proving the flat-portion
assertion without any exception at its right endpoint.
\end{proof}
\end{annexitem}
